\IfFileExists{stacks-project.cls}{%
\documentclass{stacks-project}
}{%
\documentclass{amsart}
}

% For dealing with references we use the comment environment
\usepackage{verbatim}
\newenvironment{reference}{\comment}{\endcomment}
%\newenvironment{reference}{}{}
\newenvironment{slogan}{\comment}{\endcomment}
\newenvironment{history}{\comment}{\endcomment}

% For commutative diagrams we use Xy-pic
\usepackage[all]{xy}

% We use 2cell for 2-commutative diagrams.
\xyoption{2cell}
\UseAllTwocells

% We use multicol for the list of chapters between chapters
\usepackage{multicol}

% This is generally recommended for better output
\usepackage{lmodern}
\usepackage[T1]{fontenc}

% For cross-file-references
\usepackage{xr-hyper}

% Package for hypertext links:
\usepackage{hyperref}

% For any local file, say "hello.tex" you want to link to please
% use \externaldocument[hello-]{hello}
\externaldocument[introduction-]{introduction}
\externaldocument[conventions-]{conventions}
\externaldocument[sets-]{sets}
\externaldocument[categories-]{categories}
\externaldocument[topology-]{topology}
\externaldocument[sheaves-]{sheaves}
\externaldocument[sites-]{sites}
\externaldocument[stacks-]{stacks}
\externaldocument[fields-]{fields}
\externaldocument[algebra-]{algebra}
\externaldocument[brauer-]{brauer}
\externaldocument[homology-]{homology}
\externaldocument[derived-]{derived}
\externaldocument[simplicial-]{simplicial}
\externaldocument[more-algebra-]{more-algebra}
\externaldocument[smoothing-]{smoothing}
\externaldocument[modules-]{modules}
\externaldocument[sites-modules-]{sites-modules}
\externaldocument[injectives-]{injectives}
\externaldocument[cohomology-]{cohomology}
\externaldocument[sites-cohomology-]{sites-cohomology}
\externaldocument[dga-]{dga}
\externaldocument[dpa-]{dpa}
\externaldocument[sdga-]{sdga}
\externaldocument[hypercovering-]{hypercovering}
\externaldocument[schemes-]{schemes}
\externaldocument[constructions-]{constructions}
\externaldocument[properties-]{properties}
\externaldocument[morphisms-]{morphisms}
\externaldocument[coherent-]{coherent}
\externaldocument[divisors-]{divisors}
\externaldocument[limits-]{limits}
\externaldocument[varieties-]{varieties}
\externaldocument[topologies-]{topologies}
\externaldocument[descent-]{descent}
\externaldocument[perfect-]{perfect}
\externaldocument[more-morphisms-]{more-morphisms}
\externaldocument[flat-]{flat}
\externaldocument[groupoids-]{groupoids}
\externaldocument[more-groupoids-]{more-groupoids}
\externaldocument[etale-]{etale}
\externaldocument[chow-]{chow}
\externaldocument[intersection-]{intersection}
\externaldocument[pic-]{pic}
\externaldocument[weil-]{weil}
\externaldocument[adequate-]{adequate}
\externaldocument[dualizing-]{dualizing}
\externaldocument[duality-]{duality}
\externaldocument[discriminant-]{discriminant}
\externaldocument[derham-]{derham}
\externaldocument[local-cohomology-]{local-cohomology}
\externaldocument[algebraization-]{algebraization}
\externaldocument[curves-]{curves}
\externaldocument[resolve-]{resolve}
\externaldocument[models-]{models}
\externaldocument[functors-]{functors}
\externaldocument[equiv-]{equiv}
\externaldocument[pione-]{pione}
\externaldocument[etale-cohomology-]{etale-cohomology}
\externaldocument[proetale-]{proetale}
\externaldocument[relative-cycles-]{relative-cycles}
\externaldocument[more-etale-]{more-etale}
\externaldocument[trace-]{trace}
\externaldocument[crystalline-]{crystalline}
\externaldocument[spaces-]{spaces}
\externaldocument[spaces-properties-]{spaces-properties}
\externaldocument[spaces-morphisms-]{spaces-morphisms}
\externaldocument[decent-spaces-]{decent-spaces}
\externaldocument[spaces-cohomology-]{spaces-cohomology}
\externaldocument[spaces-limits-]{spaces-limits}
\externaldocument[spaces-divisors-]{spaces-divisors}
\externaldocument[spaces-over-fields-]{spaces-over-fields}
\externaldocument[spaces-topologies-]{spaces-topologies}
\externaldocument[spaces-descent-]{spaces-descent}
\externaldocument[spaces-perfect-]{spaces-perfect}
\externaldocument[spaces-more-morphisms-]{spaces-more-morphisms}
\externaldocument[spaces-flat-]{spaces-flat}
\externaldocument[spaces-groupoids-]{spaces-groupoids}
\externaldocument[spaces-more-groupoids-]{spaces-more-groupoids}
\externaldocument[bootstrap-]{bootstrap}
\externaldocument[spaces-pushouts-]{spaces-pushouts}
\externaldocument[spaces-chow-]{spaces-chow}
\externaldocument[groupoids-quotients-]{groupoids-quotients}
\externaldocument[spaces-more-cohomology-]{spaces-more-cohomology}
\externaldocument[spaces-simplicial-]{spaces-simplicial}
\externaldocument[spaces-duality-]{spaces-duality}
\externaldocument[formal-spaces-]{formal-spaces}
\externaldocument[restricted-]{restricted}
\externaldocument[spaces-resolve-]{spaces-resolve}
\externaldocument[formal-defos-]{formal-defos}
\externaldocument[defos-]{defos}
\externaldocument[cotangent-]{cotangent}
\externaldocument[examples-defos-]{examples-defos}
\externaldocument[algebraic-]{algebraic}
\externaldocument[examples-stacks-]{examples-stacks}
\externaldocument[stacks-sheaves-]{stacks-sheaves}
\externaldocument[criteria-]{criteria}
\externaldocument[artin-]{artin}
\externaldocument[quot-]{quot}
\externaldocument[stacks-properties-]{stacks-properties}
\externaldocument[stacks-morphisms-]{stacks-morphisms}
\externaldocument[stacks-limits-]{stacks-limits}
\externaldocument[stacks-cohomology-]{stacks-cohomology}
\externaldocument[stacks-perfect-]{stacks-perfect}
\externaldocument[stacks-introduction-]{stacks-introduction}
\externaldocument[stacks-more-morphisms-]{stacks-more-morphisms}
\externaldocument[stacks-geometry-]{stacks-geometry}
\externaldocument[moduli-]{moduli}
\externaldocument[moduli-curves-]{moduli-curves}
\externaldocument[examples-]{examples}
\externaldocument[exercises-]{exercises}
\externaldocument[guide-]{guide}
\externaldocument[desirables-]{desirables}
\externaldocument[coding-]{coding}
\externaldocument[obsolete-]{obsolete}
\externaldocument[fdl-]{fdl}
\externaldocument[index-]{index}

% Theorem environments.
%
\theoremstyle{plain}
\newtheorem{theorem}[subsection]{Theorem}
\newtheorem{proposition}[subsection]{Proposition}
\newtheorem{lemma}[subsection]{Lemma}

\theoremstyle{definition}
\newtheorem{definition}[subsection]{Definition}
\newtheorem{example}[subsection]{Example}
\newtheorem{exercise}[subsection]{Exercise}
\newtheorem{situation}[subsection]{Situation}

\theoremstyle{remark}
\newtheorem{remark}[subsection]{Remark}
\newtheorem{remarks}[subsection]{Remarks}

\numberwithin{equation}{subsection}

% Macros
%
\def\lim{\mathop{\mathrm{lim}}\nolimits}
\def\colim{\mathop{\mathrm{colim}}\nolimits}
\def\Spec{\mathop{\mathrm{Spec}}}
\def\Hom{\mathop{\mathrm{Hom}}\nolimits}
\def\Ext{\mathop{\mathrm{Ext}}\nolimits}
\def\SheafHom{\mathop{\mathcal{H}\!\mathit{om}}\nolimits}
\def\SheafExt{\mathop{\mathcal{E}\!\mathit{xt}}\nolimits}
\def\Sch{\mathit{Sch}}
\def\Mor{\mathop{\mathrm{Mor}}\nolimits}
\def\Ob{\mathop{\mathrm{Ob}}\nolimits}
\def\Sh{\mathop{\mathit{Sh}}\nolimits}
\def\NL{\mathop{N\!L}\nolimits}
\def\CH{\mathop{\mathrm{CH}}\nolimits}
\def\proetale{{pro\text{-}\acute{e}tale}}
\def\etale{{\acute{e}tale}}
\def\QCoh{\mathit{QCoh}}
\def\Ker{\mathop{\mathrm{Ker}}}
\def\Im{\mathop{\mathrm{Im}}}
\def\Coker{\mathop{\mathrm{Coker}}}
\def\Coim{\mathop{\mathrm{Coim}}}

% Boxtimes
%
\DeclareMathSymbol{\boxtimes}{\mathbin}{AMSa}{"02}

%
% Macros for moduli stacks/spaces
%
\def\QCohstack{\mathcal{QC}\!\mathit{oh}}
\def\Cohstack{\mathcal{C}\!\mathit{oh}}
\def\Spacesstack{\mathcal{S}\!\mathit{paces}}
\def\Quotfunctor{\mathrm{Quot}}
\def\Hilbfunctor{\mathrm{Hilb}}
\def\Curvesstack{\mathcal{C}\!\mathit{urves}}
\def\Polarizedstack{\mathcal{P}\!\mathit{olarized}}
\def\Complexesstack{\mathcal{C}\!\mathit{omplexes}}
% \Pic is the operator that assigns to X its picard group, usage \Pic(X)
% \Picardstack_{X/B} denotes the Picard stack of X over B
% \Picardfunctor_{X/B} denotes the Picard functor of X over B
\def\Pic{\mathop{\mathrm{Pic}}\nolimits}
\def\Picardstack{\mathcal{P}\!\mathit{ic}}
\def\Picardfunctor{\mathrm{Pic}}
\def\Deformationcategory{\mathcal{D}\!\mathit{ef}}
\usepackage{cleveref}
\crefname{originalreference}{original reference}{original references}
\Crefname{originalreference}{Original reference}{Original references}
\makeatletter
\newlabel{chow-situation-setup-base-change}{{0FVG}{source}{}{sourcepointer-1}{}}
\newlabel{chow-situation-setup-base-change@cref}{{[originalreference][0][]0FVG}{[1][0][]source}{}{}{}}
\makeatother
\crefname{originalsection}{original section}{original sections}
\Crefname{originalsection}{Original section}{Original sections}
\makeatletter
\newlabel{chow-section-change-base}{{0FVF}{source}{}{sourcepointer-2}{}}
\newlabel{chow-section-change-base@cref}{{[originalsection][0][]0FVF}{[1][0][]source}{}{}{}}
\makeatother
\crefname{originallemma}{original lemma}{original lemmas}
\Crefname{originallemma}{Original lemma}{Original lemmas}
\makeatletter
\newlabel{lemma-action-opens}{{0H67}{source}{}{sourcepointer-3}{}}
\newlabel{lemma-action-opens@cref}{{[originallemma][0][]0H67}{[1][0][]source}{}{}{}}
\makeatother
\crefname{originallemma}{original lemma}{original lemmas}
\Crefname{originallemma}{Original lemma}{Original lemmas}
\makeatletter
\newlabel{lemma-action-bilinear}{{0H66}{source}{}{sourcepointer-4}{}}
\newlabel{lemma-action-bilinear@cref}{{[originallemma][0][]0H66}{[1][0][]source}{}{}{}}
\makeatother
\crefname{originallemma}{original lemma}{original lemmas}
\Crefname{originallemma}{Original lemma}{Original lemmas}
\makeatletter
\newlabel{chow-lemma-pullback-base-change-pullback}{{0FVK}{source}{}{sourcepointer-5}{}}
\newlabel{chow-lemma-pullback-base-change-pullback@cref}{{[originallemma][0][]0FVK}{[1][0][]source}{}{}{}}
\makeatother
\title{Compatibility of the cycle intersection with flat base change and flat pullback}
\author{The Stacks Project Authors}
\date{stacks-project:relative-cycles.tex:local-snapshot-20261012; GNU FDL 1.2 or later}
\begin{document}
\maketitle
\setcounter{subsection}{1}
\noindent\textit{Editorial note.} All delivered mathematical text is the author's own, copied from the pinned Stacks Project chapter file relative-cycles.tex; no statement or proof step was rewritten, reordered or removed. The counted span is the author environment and its own proof environment. References inside the span are handed to the assembler as navigation stubs that carry the canonical tag of the cited tag only; whether each of them is load-bearing is recorded on the stub itself.\par\medskip
\noindent\textit{Allowed standard prerequisites (editorial).}\begin{itemize}
\item Chow Homology Situation setting up base change of cycles with a morphism g : S' to S and a function c on the delta dimensions (chow chapter, Situation chow-situation-setup-base-change, canonical tag 0FVG). cited at proof line 2265. Fixes the base-change setup: delta' = delta + c and the meaning of the base change of a family of cycles.
\item Chow Homology section on change of base for cycles, in particular that the pullback of a prime cycle beta is the sum over irreducible components with multiplicities (chow chapter, Section chow-section-change-base, canonical tag 0FVF). cited at proof line 2272. Source of the formula beta' = sum n\_a [Z'\_a] used in the multiplicity computation.
\item Lemma allowing the intersection operation to be checked locally on Y: for open U in X and V in Y with f(U) in V, restricting the factors restricts the intersection (more precisely the relative-cycles chapter, Lemma lemma-action-opens, canonical tag 0H67; rejected as a standalone card, cited here as a named input). cited at proof line 2279. Licences the reduction to Y affine so that beta is a finite linear combination of prime cycles.
\item Bilinearity of the intersection operation in the second variable (relative-cycles chapter, Lemma lemma-action-bilinear, canonical tag 0H66; rejected as a standalone card, cited here as a named input). cited at proof line 2282. Reduces the equality from arbitrary beta to the prime-cycle case beta = [Z].
\item Commutativity of the square relating base change of zero-cycles on the generic fibre and flat pullback (chow chapter, Lemma chow-lemma-pullback-base-change-pullback, canonical tag 0FVK). cited at proof line 2344. The named endpoint: the diagram on cycle groups base-changed by the generic point and pulled back flat is commutative, which is exactly the component and multiplicity claim after the birational transport.
\end{itemize}\medskip
\ifdefined\hypertarget\hypertarget{sourcepointer-1}{}\fi
\noindent\textit{Source reference.} Original reference 0FVG: sources/stacks-chow/tags.txt, source lines 19811--19811. Cited at proof line 2265: fixes the base-change setup (delta prime equals delta plus c) for the whole proof.\par\medskip
\ifdefined\hypertarget\hypertarget{sourcepointer-2}{}\fi
\noindent\textit{Source reference.} Original section 0FVF: sources/stacks-chow/tags.txt, source lines 19810--19810. Cited at proof line 2272: source of the pullback formula for beta with component multiplicities.\par\medskip
\ifdefined\hypertarget\hypertarget{sourcepointer-3}{}\fi
\noindent\textit{Source reference.} Original lemma 0H67: sources/stacks-chow/tags.txt, source lines 21412--21412. Cited at proof line 2279: licences working locally on Y so beta splits into prime cycles.\par\medskip
\ifdefined\hypertarget\hypertarget{sourcepointer-4}{}\fi
\noindent\textit{Source reference.} Original lemma 0H66: sources/stacks-chow/tags.txt, source lines 21411--21411. Cited at proof line 2282: reduces the equality to the prime-cycle case.\par\medskip
\ifdefined\hypertarget\hypertarget{sourcepointer-5}{}\fi
\noindent\textit{Source reference.} Original lemma 0FVK: sources/stacks-chow/tags.txt, source lines 19815--19815. Cited at proof line 2344: the named commutative square that finishes the multiplicity computation.\par\medskip
\begin{lemma}
\label{lemma-action-base-change}
Forming $\alpha \cap \beta$ is compatible with flat base change
and flat pullback (see proof for elucidation).
\end{lemma}

\begin{proof}
Let $(S, \delta)$, $(S', \delta')$, $g : S' \to S$, and $c \in \mathbf{Z}$
be as in Chow Homology, Situation \ref{chow-situation-setup-base-change}.
Let $X \to Y$ be a morphism of schemes locally of finite type over $S$.
Denote $X' \to Y'$ the base change of $X \to Y$ by $g$.
Let $\alpha$ be a family of $r$-cycles on the fibres of $X/Y$.
Let $\beta \in Z_e(Y)$. Denote $\alpha'$ the
base change of $\alpha$ by $Y' \to Y$. Denote
$\beta' = g^*\beta \in Z_{e + c}(Y')$ the pullback of $\beta$ by $g$, see
Chow Homology, Section \ref{chow-section-change-base}.
Compatibility with base change means
$\alpha' \cap \beta'$ is the base change of $\alpha \cap \beta$.

\medskip\noindent
Proof of compatibility with base change.
Since we are proving an equality of cycles on $X'$, we may work locally
on $Y$, see Lemma \ref{lemma-action-opens}. Thus we may assume $Y$
is affine. In particular $\beta$ is a finite linear combination of prime
cycles. Since $- \cap -$ is linear in the second variable
(Lemma \ref{lemma-action-bilinear}), it suffices
to prove the equality when $\beta = [Z]$ for some integral closed subscheme
$Z \subset Y$ of $\delta$-dimension $e$.

\medskip\noindent
Let $y \in Z$ be the generic point. Write $\alpha_y = \sum m_j [V_j]$.
Let $\overline{V}_j$ be the closure of $V_j$ in $X$. Then we have
$$
\alpha \cap \beta = \sum m_j[\overline{V}_j]
$$
The base change of $\beta$ is $\beta' = \sum [Z \times_S S']_{e + c}$ as a
cycle on $Y' = Y \times_S S'$. Let $Z'_a \subset Z \times_S S'$
be the irreducible components, denote $y'_a \in Z'_a$ their generic points,
and denote $n_a$ the multiplicity of $Z'_a$ in $Z \times_S S'$.
We have
$$
\beta' = \sum [Z \times_S S']_{e + c} = \sum n_a[Z'_a]
$$
We have $\alpha'_{y'_a} = \sum m_j [V_{j, \kappa(y'_a)}]_r$
because $\alpha'$ is the base change of $\alpha$ by $Y' \to Y$.
Let $V'_{jab} \subset V_{j, \kappa(y'_a)}$ be the irreducible components
and denote $m_{jab}$ the multiplicity of $V'_{jab}$
in $V_{j, \kappa(y'_a)}$. We have
$$
\alpha'_{y'_a} = \sum m_j [V_{j, \kappa(y'_a)}]_r =
\sum m_j m_{jab} [V'_{jab}]
$$
Thus we we have
$$
\alpha' \cap \beta' = \sum n_a m_j m_{jab} [\overline{V}'_{jab}]
$$
where $\overline{V}'_{jab}$ is the closure of $V'_{jab}$ in $X'$.
Thus to prove the desired equality it suffices to prove
\begin{enumerate}
\item the irreducible components of $\overline{V}_j \times_S S'$
are the schemes $\overline{V}'_{jab}$ and
\item the multiplicity of $\overline{V}'_{jab}$ in
$\overline{V}_j \times_S S'$ is equal to $n_a m_{jab}$.
\end{enumerate}
Note that $V_j \to \overline{V}_j$ is a birational morphism of
integral schemes. The morphisms $V_j \times_S S' \to V_j$
and $\overline{V}_j \times_S S' \to \overline{V}_j$ are flat
and hence map generic points of irreducible components to the
(unique) generic points of $V_j$ and $\overline{V}_j$.
It follows that $V_j \times_S S' \to \overline{V}_j \times_S S'$
is a birational morphisms hence induces a bijection on irreducible
components and identifies their multiplicities.
This means that it suffices to prove that the irreducible components of
$V_j \times_S S'$ are the schemes $V'_{jab}$ and
the multiplicity of $V'_{jab}$ in $V_j \times_S S'$
is equal to $n_a m_{jab}$. However, then we are just saying
that the diagram
$$
\xymatrix{
Z_r(V_j) \ar[r] & Z_{r + c}(V_j \times_S S') \\
Z_0(\Spec(\kappa(y))) \ar[r] \ar[u] &
Z_c(\Spec(\kappa(y)) \times_S S') \ar[u]
}
$$
is commutative where the horizontal arrows are base change by
$\Spec(\kappa(y)) \times_S S' \to \Spec(\kappa(y))$ and
the vertical arrows are flat pullback. This was shown in
Chow Homology, Lemma \ref{chow-lemma-pullback-base-change-pullback}.

\medskip\noindent
The statement in the lemma on flat pullback means the following.
Let $(S, \delta)$, $X \to Y$, $\alpha$, and $\beta$ be as in the
construction of $\alpha \cap \beta$ above. Let $Y' \to Y$ be a flat
morphism, locally of finite type, and of relative dimension $c$. Then
we can let $\alpha'$ be the base change of $\alpha$ by $Y' \to Y$
and $\beta'$ the flat pullback of $\beta$. Compatibility with flat
pullback means $\alpha' \cap \beta'$ is the flat pullback of
$\alpha \cap \beta$ by $X \times_Y Y' \to Y$. This is actually
a special case of the discussion above if we set $S = Y$ and $S' = Y'$.
\end{proof}

\end{document}
